%--------------------------------------------------------------------------------
%
%
%--------------------------------------------------------------------------------
\pagebreak
\section*{HELP - list help info}
\addcontentsline{toc}{section}{HELP - list help info}

%--------------------------------------------------------------------------------
\begin{iPageDescEntry}{Syntax}
\texttt{HELP } \\
\texttt{HELP <cmdId> } \\
\texttt{HELP "commands" } \\
\texttt{HELP "wcommands" } \\ 
\texttt{HELP "predefined" }
\end{iPageDescEntry}

%--------------------------------------------------------------------------------
\begin{iPageDescEntry}{Description}

The \texttt{HELP} command displays help information for the simulator. Help topics are grouped into several categories:

\begin{itemize}
  \item The \texttt{COMMANDS} keyword lists a summary of all available line-mode commands. Entering \texttt{HELP} with no arguments also displays this list.
  \item The \texttt{WCOMMANDS} keyword lists a summary of all available window-mode commands.
  \item The \texttt{PREDEFINED} keyword lists all predefined functions that can be used in command-line expressions.
\end{itemize}

To view detailed syntax for a specific command or predefined function, type \texttt{HELP <cmdId>}
where \textit{cmdId} is the name of the command or function.


\end{iPageDescEntry}

%--------------------------------------------------------------------------------
\begin{iPageDescOpEntry}{Example}
\begin{lstlisting}[style=iPageOpStyle]
(8) ->help
 help            list help info
 exit            program exit
 hist            command history
 do              re-execute command
 redo            edit and then re-execute command
 env             lists the env tab, a variable, sets a variable
 xf              execute commands from a file
 reset           resets the system, etc.
 run             run the system ( all CPUs )
 step            single step the system
 w               evaluates and prints an expression
 dm              display registered modules for the system
 mr              modify a CPU register
 da              display absolute memory
 ma              modify absolute memory
 pica            purges instruction cache line data
 pdca            purges data cache line data
 fdca            flushes data cache line data
 iitlb           inserts an entry into the instruction TLB
 idtlb           inserts an entry into the data TLB
 pitlb           purges an entry from the instruction TLB
 pdtlb           purges an entry from the data TLB
 won             enables windows mode / redraw
(9) ->
(9) ->help da
da <adr> [ , <len> ] [ , <fmt> ] - display absolute memory
(10) ->
\end{lstlisting}
\end{iPageDescOpEntry}

%--------------------------------------------------------------------------------
\begin{iPageDescEntry}{Notes}
None.
\end{iPageDescEntry}
